\begin{figure}[htbp]
\centering
\begin{tikzpicture}[
    >=Latex,
    font=\small,
    box/.style={draw, semithick, rounded corners=2pt, align=center,
                minimum height=10mm, minimum width=39mm, inner sep=4pt, fill=black!2},
    key/.style={box, fill=black!7},
    arr/.style={->, semithick}
]
\node[key] (input) at (0,4.4) {High-density SNP\\genotype matrix};
\node[box] (gen) at (-3.3,2.9) {Variational\\model};
\node[box] (con) at (3.3,2.9) {Contrastive\\model};
\node[key] (repr) at (0,1.45) {Learned genomic\\representation};
\node[box] (cls) at (-3.3,0.0) {Breed\\classification};
\node[box] (xai) at (3.3,0.0) {Feature\\attribution};
\node[key] (rank) at (3.3,-1.45) {Global SNP\\ranking};
\node[key] (panel) at (3.3,-2.9) {Reduced SNP\\panels};
\node[box] (sample) at (-3.3,-2.9) {Sample-size\\analysis};
\node[box] (marker) at (3.3,-4.35) {Marker-count\\analysis};
\node[key] (cross) at (0,-5.8) {Cross-species validation\\caprine $\rightarrow$ ovine};

\draw[arr] (input) -- (gen);
\draw[arr] (input) -- (con);
\draw[arr] (gen) -- (repr);
\draw[arr] (con) -- (repr);
\draw[arr] (repr) -- (cls);
\draw[arr] (repr) -- (xai);
\draw[arr] (xai) -- (rank);
\draw[arr] (rank) -- (panel);
\draw[arr] (cls) -- (sample);
\draw[arr] (panel) -- (marker);
\draw[arr] (sample) -- (cross);
\draw[arr] (marker) -- (cross);
\end{tikzpicture}

\caption{High-level computational organization of the thesis. The workflow links representation learning to breed classification and SNP attribution, then evaluates the effects of training-set size and marker-panel size before cross-species validation. Author's elaboration.}
\label{fig:thesis_overview}
\end{figure}
